\documentclass{article}

\usepackage{amssymb,amsmath,amsthm,mathtools,cancel}
\usepackage{tikz}
\usepackage[margin=1in]{geometry}
\usepackage{enumitem}

\newtheorem{theorem}{Theorem}[section]
\newtheorem{lemma}[theorem]{Lemma}

\begin{document}
\noindent
\begin{minipage}[t]{0.68\textwidth}
{\large 6.1200 \textit{Mathematics for Computer Science}}\\[4pt]
{\large Problem Set 8}\\[4pt]
{\large Solutions by Harsh Bhalala}
\end{minipage}%
\hfill
\begin{minipage}[t]{0.27\textwidth}
\raggedleft
{\large 8, October 2026}
\end{minipage}

\vspace{6pt}

\noindent\rule{\textwidth}{0.5pt}

\section*{Problem 1. Strong Connectivity}

\textbf{(a)} In a directed graph $G$ equivalence relation $R$ on the vertices of $G$ as follows: $u \mathrel{R} v$
is true IFF there exists a directed walk from $u$ to $v$ and a directed walk from $v$ to $u$.
(Note: A walk from $ u$ to $ v$ is expressed as $u \rightarrow v$)
Now to show that the given relation is an equivalence relation:
\begin{itemize}
\item \textbf{Reflexive:} For any vertex u, there always exists walk of length from $u$ to $u$ in both directions, so $u \mathrel{R} u$.
\item \textbf{Symmetric :} Suppose $u \mathrel{R} v$, There then there exists $(u \rightarrow v)$ and $(v \rightarrow u)$,Which are exactly two
walks required for relation $u \mathrel{R} v$.
\item \textbf{Transitive:} If $a \mathrel{R} b$ and $b \mathrel{R} c$ then we know there exists walks $(a \rightarrow b)$, $(b \rightarrow a)$,
$(b \rightarrow c)$, $(c \rightarrow b)$ and by concatenating those walk we get $(a \rightarrow b \rightarrow c)$ which imples $(a \rightarrow c)$
and similarly, by doing the same for the other two walks we get ($c \rightarrow a$). Hence $b \mathrel{R} c$.
\end{itemize}

\vspace{\baselineskip}
\noindent\textbf{(b)}
By comparing the definition of strongly connected components and $R$ we observe that The equivalence classes of $R$ are exactly the strongly connected components.

\section*{Problem 2. Bijections via Inverses}

\textbf{(a)}
\begin{proof}[Proof]
Let $a_1,a_2 \in A$ such that $f(a_1) = f(a_2)$, now applying $g$ to both side, $g(f(a_1)) = g(f(a_2))$ since $\forall a, g(f(a)) = a$,
$a_1 = a_2$. So $f$ is surjective. 
\end{proof}

\vspace{\baselineskip}
\noindent\textbf{(b)}
\begin{proof}[proof by contradiction]
Assume $f$ is not surjective. Hence, $\exists b \in B$ such that $\forall a \in A, f(a) \neq b$.But $f(g(b)) = b$ is true for all $b\in B$
which is a contradiction since $g(b)$ is an element of A that maps to b.
\end{proof}

\noindent\textbf{(c)} From parts (a) and (b), we know $f$ is a function, total, injective, and surjective. Hence $f$ is a bijection.


\section*{Problem 3. Increasing Correspondence}

\textbf{(a)}
Let $(a,b,c) \in A$. Now by the definition $f(a,b,c) = (c-b,b-a,a)$ and because $a \leq b \leq c$ we observe that $c-b \geq 0$,
$b-a \geq 0$ and $a \geq 0$, which implies that for every valid input pair there is at least one output.
Therefore, the given function is a total function.

Now we show that every output of $f$ lies in set B. As we know $\forall(a,b,c) \in A$, $f(a,b,c) = ((c-b,b-a,a))$, and given pair is
part of set B IFF
\[
\begin{aligned}
(c-b) + 2(b-a) + 3a = a + b + c = 1000\\
\end{aligned}
\]
which satisfies all necessary conditions and by this we can conclude that $\forall (a, b, c) \in A$,
so $f((a, b, c)) \in B$.

\pagebreak

\vspace{\baselineskip}
\noindent\textbf{(b)}
Total function $g:B \rightarrow A$, $g(x,y,z) \coloneq (z, y+z, x+y+z)$

$\forall (x,y,z) \in B$, triples (z,y,x) is defined over natural number and since addition is deterministic operation (z,y+z,x+y+z) will
garentes exactly one output for all possible input. Also $x+y+z \geq y+z \geq z$ is true because $y,x \geq 0$. 
Sum: $z + (y+z) + (x+y+z) = x + 2y + 3z = 1000$. So $g(x,y,z) \in A$.
\[
\begin{aligned}
    g \circ f &= g(f(a,b,c)) = g(c-b, b-a, a) = (a, (b-a)+a, (c-b)+(b-a)+a) = (a, b, c).\\ 
    f \circ  g &= f(g(x,y,z)) = f(z, y+z, x+y+z) = ((x+y+z)-(y+z), (y+z)-z, z) = (x, y, z).\\
\end{aligned}
\]
\vspace{\baselineskip}
\noindent\textbf{(c)} Since we proved that $g$ is a two-sided inverse of $f$, and that is necessary and sufficient to prove that $f$
is bijection


\section*{Problem 4. Counting Graphs}
\textbf{(a)} \[2^{\frac{n(n-1)}{2}}\]


\vspace{\baselineskip}
\noindent\textbf{(b)} \[n \times 4 \times 3 \times 3 = 36n\]



\section*{Problem 5. Counting}
\textbf{(a)} There are toal 21! ways to arrange 21 consonants
Now, since every vowel appears to the left of a consonant,t we must choose 5 total consonants,nt so make a gap in front and place a vowel
. There are $\binom{21}{5}$ ways to make such a gap, and now there are 5! ways to put a vowel into those gaps.

Hence total way = \[21! \times \binom{21}{5} \times 5!\]

\vspace{\baselineskip}
\noindent\textbf{(b)}
We first need to make a total of 5 [VCC] boxes and the remaining 11 consonants in [C] boxes:
\[ \frac{\left( \binom{21}{2} \times \binom{19}{2} \times \binom{17}{2} \times \binom{15}{2} \times \binom{13}{2} \times 2^5 \right)}{5!} \times 5! \times 16!\]
We needed to divide by 5! because the count is for distinct sections of groups, meaning (bc, df,...) is the same as (df, bc,...)
and $2^5$ because pairs bc and cb are distinct. and there are 5! ways to place vowels in each box.
and there are a total of 11 boxes of [C], so the number of ways to arrange 16 distinct boxes is 16!

And simplifying the given equation gives:
\[\frac{21! \times 16!}{11!}\]


\vspace{\baselineskip}
\noindent\textbf{(c)} (Note: Another concise way to expressthe  equation in (b) is used here, which is viathe  multinomial coefficient)
Need to choose $ n$ partitions each of size 2, from 2n choises that can be done in:
\[\binom{2n}{2,2,2,\dots 2} \times \frac{1}{n!} = \frac{(2n)!}{2^n n!}\] We must here, just like in part (b), the binomial coefficient gives an ordered pair, and since there are a total of $ n$ terms, there are $ n! $ copies of
same arrangement.

\vspace{\baselineskip}
\noindent\textbf{(d)}We can simply apply the Stars and Bars combinatorics theorem,Here we have n stars and 9 bars,
(since there are 10 digits 0 to 9 and each partition says how many times a digit appears in the number) 
and we need to find the number of ways to arrange them. 
The total number of positions to place bars is $n + 9$ and we need to choose 9 positions for bars, which gives us the following expression:
\[ \binom{\text{total position to place bars}}{\text{bars}} = \binom{\text{stars+bars}}{\text{bars}} = \binom{n + 9}{9}\]

\end{document}